% TOP_PROBLEM: {"id":30007634,"problem_number":"LOCAL-30007634","title":"Job-search information and persistent mismatch","category_id":21,"category":{"id":21,"name":"economics","display_name":"Economics","description":"Open economic theory statements, published research agendas, and proposed scoped specifications; question provenance and historical priority are qualified per record.","slug":"economics","order_index":21,"created_at":"2026-10-08T00:00:00Z"},"status":"research_question","external_url":null,"legacy_ids":[],"published":false,"statement":"In the explicitly specified setting: How much inefficient job matching comes from mistaken beliefs about pay and employers, and which information interventions durably improve matches at scale? The mathematical statement above fixes the target, assumptions and scope of this catalogue formulation.","economics":{"original_id":123,"field":"Labor markets and work","kind":"identification","formalization_basis":"adapted_research_question","source_status":"research_agenda","priority_status":"earliest_found","priority_note":"Earliest directly inspected qualifying passage in this audit; earlier origins were not established. A review or research-agenda statement does not imply original priority.","earliest_verified_reference":{"authors":["Sydnee Caldwell","Ingrid Haegele","Jörg Heining"],"title":"Why Workers Stay: Pay, Beliefs, and Attachment","year":2025,"url":"https://sydneec.github.io/Website/CHH_Monopsony.pdf","doi":"10.3386/w33445","locator":"December 2025 version, conclusion, printed p. 41, final paragraph","evidence_summary":"The conclusion calls for work on social ties and uncertainty about particular nonpay attributes; pay misinformation alone is not their dominant explanation."},"current_reference":{"authors":["Sydnee Caldwell","Ingrid Haegele","Jörg Heining"],"title":"Why Workers Stay: Pay, Beliefs, and Attachment","year":2025,"url":"https://sydneec.github.io/Website/CHH_Monopsony.pdf","doi":"10.3386/w33445","locator":"December 2025 version, conclusion, printed p. 41, final paragraph","evidence_summary":"The conclusion calls for work on social ties and uncertainty about particular nonpay attributes; pay misinformation alone is not their dominant explanation."},"formalization_note":"The cited passage establishes a research direction, not this exact mathematical statement. The notation, population, policy menu, assumptions, and estimands are an original adaptation for this catalogue; no universal unsolved theorem or source priority is claimed."},"statusQualification":"Published question or research agenda verified at the cited locator. The mathematical specification is adapted for this catalogue and includes additional assumptions; source publication does not establish first-ever priority or certify this exact model remains unsolved. The cited passage establishes a research direction, not this exact mathematical statement. The notation, population, policy menu, assumptions, and estimands are an original adaptation for this catalogue; no universal unsolved theorem or source priority is claimed.","statusReviewedAt":"2026-10-08"}
% FIRST_POSTED: 2026-10-08
% ENTRY_KIND: statement_only
% !TeX program = lualatex
\documentclass[11pt]{article}
\usepackage{fontspec}
\usepackage[a4paper,margin=25mm]{geometry}
\usepackage{amsmath,amssymb,mathtools}
\usepackage{xurl}
\usepackage[hidelinks]{hyperref}
\newcommand{\sref}[2]{\hyperlink{source:#1:#2}{[S#2]}}
\setlength{\emergencystretch}{3em}
\begin{document}
% problemId:30007634
\section*{Job-search information and persistent mismatch}\label{30007634}
\subsection{Definitions and mathematical statement}
Consider workers searching within a specified labor market. For worker \(i\) and employer \(j\), let \(b_{ij}\) be the worker's belief about pay and nonpay job characteristics, \(v_{ij}\) the value of a fully informed feasible offer, and \(k_{ij}\) the switching cost. An information intervention \(z\in\{0,1\}\) supplies independently validated employer information; it does not guarantee an offer. Let \(J_i(z)\) be the eventual employer and \(Y_{iT}(z)\) discounted earnings through a fixed horizon \(T\). Define \(\tau_T=E[Y_{iT}(1)-Y_{iT}(0)]\) and \(\tau_J=\Pr(J_i(1)\ne J_i(0))\). Randomization identifies marginal employer distributions and \(\tau_T\); the cross-world switching rate \(\tau_J\) requires coupling assumptions or sharp bounds. The task is to distinguish corrections of mistaken beliefs from changes in preferences, social ties, or offer probabilities. A complementary structural target is the loss in expected job value attributable to inaccurate \(b_{ij}\), relative to truthful information with the same feasible offers. Random assignment identifies information effects only with consistency, appropriate interference restrictions, and measured attrition. At scale, employer responses and vacancy competition must be included. Determine when the intervention improves durable earnings and worker value, and characterize which belief dimensions and worker groups account for heterogeneity.

\par\textbf{Formulation and scope.} The cited passage establishes a research direction, not this exact mathematical statement. The notation, population, policy menu, assumptions, and estimands are an original adaptation for this catalogue; no universal unsolved theorem or source priority is claimed.

\par\textbf{Open-question provenance.} An explicit question or research agenda was verified at the source locator. This does not by itself certify that every later version remains unresolved \sref{30007634}{1}.

\par\textbf{Historical priority.} Earliest directly inspected qualifying passage in this audit; earlier origins were not established. A review or research-agenda statement does not imply original priority.

\subsection{Short English statement}
In the explicitly specified setting: How much inefficient job matching comes from mistaken beliefs about pay and employers, and which information interventions durably improve matches at scale? The mathematical statement above fixes the target, assumptions and scope of this catalogue formulation.

\subsection{Sources}
\begin{itemize}\raggedright
\item[\textnormal{[S1]}]\hypertarget{source:30007634:1}{} Sydnee Caldwell, Ingrid Haegele, Jörg Heining. 2025. Why Workers Stay: Pay, Beliefs, and Attachment. December 2025 version, conclusion, printed p. 41, final paragraph.
\url{https://sydneec.github.io/Website/CHH_Monopsony.pdf}.
DOI: 10.3386/w33445.
{\small\textit{Earliest verified question/agenda reference; historical priority qualified below. The conclusion calls for work on social ties and uncertainty about particular nonpay attributes; pay misinformation alone is not their dominant explanation.}}
\end{itemize}
\end{document}
